\chapter{非负函数的积分}
\label{chap:nonnegative-integral}

Riemann 积分从定义域的分割出发，先处理区间上的阶梯状近似；Lebesgue
积分则从函数的取值出发，先计算只取有限多个值的函数，再让这些函数从下方
逼近一般非负可测函数。非负性使这一过程不需要预先讨论正负抵消，积分值也
可以是 \(+\infty\)。这一步建立以后，单调极限与积分能够交换，随后的一般
可积函数理论将由正部、负部分解得到。

\section{非负简单函数的积分}

设 \((X,\mathcal A,\mu)\) 为测度空间。非负简单函数可写成
\[
 \varphi=\sum_{k=1}^{m}a_k\mathbf 1_{E_k},
 \qquad a_k\geq 0,\quad E_k\in\mathcal A,
\]
其中 \(E_k\) 不必互不相交。沿用上一章的规范表示，在计算积分时省去零水平集，得到
\[
 \varphi=\sum_{j=1}^{r}b_j\mathbf 1_{A_j},
 \qquad b_j>0,\quad A_i\cap A_j=\varnothing\ (i\ne j).
\]

\begin{definition}
非负简单函数 \(\varphi\) 的积分定义为
\[
 \int_X\varphi\,d\mu=\sum_{j=1}^{r}b_j\mu(A_j)\in[0,+\infty],
\]
这里使用规范表示，并约定 \(0\cdot(+\infty)=0\)。对 \(E\in\mathcal A\)，
\[
 \int_E\varphi\,d\mu:=\int_X\varphi\mathbf1_E\,d\mu .
\]
\end{definition}

规范表示的作用只是便于下定义，积分并不依赖表示方法。

\begin{proposition}[良定义性与线性]
\label{prop:simple-integral-well-defined}
若
\[
 \varphi=\sum_{k=1}^{m}a_k\mathbf1_{E_k}
\]
是任意非负表示，则
\[
 \int_X\varphi\,d\mu=\sum_{k=1}^{m}a_k\mu(E_k).
\]
因此，对非负简单函数 \(\varphi,\psi\) 和 \(\alpha,\beta\ge0\)，
\[
 \int_X(\alpha\varphi+\beta\psi)\,d\mu
 =\alpha\int_X\varphi\,d\mu+\beta\int_X\psi\,d\mu .
\]
\end{proposition}

\begin{proof}
把所有集合 \(E_k\) 按“属于哪些 \(E_k\)”细分为有限个两两不交的原子。
在每个原子上，右端各项贡献之和等于 \(\varphi\) 在该原子上的常值乘以该
原子的测度。对这些原子求和即得第一式。对 \(\varphi,\psi\) 同时作这样的
有限细分，线性公式便归结为有限和的分配律。含 \(+\infty\) 时两边均按
非负扩充实数的运算法则解释，不发生 \(\infty-\infty\)。
\end{proof}

\begin{proposition}[序性质]
\label{prop:simple-integral-order}
若 \(0\le\varphi\le\psi\) 为简单函数，则
\[
 \int_E\varphi\,d\mu\le\int_E\psi\,d\mu\qquad(E\in\mathcal A).
\]
若 \(E_n\uparrow E\)，则
\[
 \int_{E_n}\varphi\,d\mu\uparrow\int_E\varphi\,d\mu .
\]
\end{proposition}

\begin{proof}
同时细分 \(\varphi,\psi\) 的水平集后，第一式逐块成立。第二式由
\(\mu(E_n\cap A_j)\uparrow\mu(E\cap A_j)\) 以及有限求和得到。
\end{proof}

\section{非负可测函数的积分}

\begin{definition}
若 \(f:X\to[0,+\infty]\) 可测，定义
\[
 \int_X f\,d\mu
 :=\sup\left\{\int_X\varphi\,d\mu:
       \varphi\ \text{为非负简单函数且 }\varphi\le f\right\}.
\]
积分允许取 \(+\infty\)。若积分有限，称 \(f\) 的非负积分有限。
对任意 \(E\in\mathcal A\)，定义
\[
 \int_E f\,d\mu=\int_X f\mathbf1_E\,d\mu,
\]
其中 \((+\infty)\cdot0=0\)，即被积函数在 \(X\setminus E\) 上取零。
若 \(f\) 本身为简单函数，序性质说明上述上确界等于原来的简单函数积分，
所以两种局部积分定义也一致。
\end{definition}

这一上确界定义不依赖某个事先选定的逼近序列。它还立即给出单调性：
\[
 0\le f\le g\quad\Longrightarrow\quad
 \int f\,d\mu\le\int g\,d\mu .
\]
对 \(c>0\)，映射 \(\varphi\mapsto c\varphi\) 将 \(f\) 下方的非负简单函数
双射到 \(cf\) 下方的非负简单函数，由简单积分的齐次性可得
\(\int cf\,d\mu=c\int f\,d\mu\)。当 \(c=0\) 时，两端按约定均为零。
加法公式需要同时
协调两个函数的简单逼近，不能只从两个彼此独立的上确界直接读出；它将在
单调收敛定理之后推出。

这里还可直接证明零测不变性，无须先用收敛定理。设非负可测函数
\(f,g\) 在一个可测零集 \(N\) 外相等。对任意简单函数 \(0\le\varphi\le f\)，
有 \(\varphi\mathbf1_{N^c}\le g\)，且简单积分满足
\(\int\varphi=\int\varphi\mathbf1_{N^c}\)。取上确界给
\(\int f\le\int g\)，交换二者即得等号。
这只说明已经可测的函数之间的积分不变，并没有断言零测集上的任意修改都可测。

\begin{example}
对 \(E\in\mathcal A\)，
\[
 \int_X\mathbf1_E\,d\mu=\mu(E).
\]
因此测度可由指示函数的积分完全恢复。计数测度下，若
\(X=\mathbb N\)，则
\[
 \int_{\mathbb N} f\,d\#=\sum_{n=0}^{\infty}f(n)
\]
对每个非负函数成立；右侧按非负级数的部分和上确界理解。
确实，简单函数的积分就是逐点求和，故每个简单下界的积分不超过右侧。
反过来，在有限集 \(\{0,\ldots,N\}\) 上取 \(\min\{f,M\}\)，集外取零，
得到一个简单下界。先令 \(M\to\infty\)，再令 \(N\to\infty\)，就取得
右侧的上确界，包括某一项或整个级数为无穷的情形。
\end{example}

\section{从下方作简单函数逼近}

\begin{theorem}[非负可测函数的简单逼近]
\label{thm:nonnegative-simple-approximation}
若 \(f:X\to[0,+\infty]\) 可测，则存在非负简单函数列
\((\varphi_n)\)，使得
\[
 0\le\varphi_n\le\varphi_{n+1}\le f,\qquad
 \varphi_n(x)\longrightarrow f(x)
\]
对每个 \(x\in X\) 成立。可以同时使
\[
 0\le f(x)-\varphi_n(x)<2^{-n}
\]
在 \(f(x)<2^n\) 的点成立。
\end{theorem}

\begin{proof}
沿用\chapterref{chap:measurable-functions}中已经验证的二进分割，
此处选取截断高度 \(2^n\) 以保留明确的误差估计。令
\[
 \varphi_n(x)=
 \begin{cases}
 2^{-n}\lfloor 2^nf(x)\rfloor,& f(x)<2^n,\\
 2^n,& f(x)\ge2^n.
 \end{cases}
\]
它只取 \(0,2^{-n},\ldots,2^n\) 这些有限多个值，各水平集由 \(f\) 的
可测水平集组成，故 \(\varphi_n\) 可测。二进网格的嵌套性给出
\(\varphi_n\le\varphi_{n+1}\)，而截断和向下取整给出
\(\varphi_n\le f\)。若 \(f(x)<\infty\)，当 \(2^n>f(x)\) 时误差小于
\(2^{-n}\)；若 \(f(x)=+\infty\)，则 \(\varphi_n(x)=2^n\to+\infty\)。
\end{proof}

由定义和该定理立刻可得：若 \(\varphi_n\uparrow f\) 是上述特定逼近，
则 \(\int f\) 至少是 \(\sup_n\int\varphi_n\)。反向不等式需要说明任意
简单函数下逼近最终都能被这列函数在积分意义上捕获，这正是单调收敛证明中
的关键一步。

\section{单调收敛定理}

\begin{theorem}[Beppo Levi 单调收敛定理]
\label{thm:monotone-convergence}
设 \(f_n,f:X\to[0,+\infty]\) 均可测，且 \(f_n\uparrow f\) 几乎处处。则
\[
 \int_X f\,d\mu=\lim_{n\to\infty}\int_X f_n\,d\mu .
\]
\end{theorem}

\begin{proof}
取一个可测零集 \(N\)，包含单调性或收敛性失效的全部点。
把 \(f\) 和所有 \(f_n\) 在整个 \(N\) 上统一置零，所得函数仍可测，
并由刚才证明的零测不变性保持积分。以下可设处处单调收敛。

由 \(f_n\le f\)，
\[
 L:=\lim_n\int f_n\,d\mu\le\int f\,d\mu .
\]
为证反向不等式，任取非负简单函数 \(\varphi\le f\) 和 \(0<c<1\)，令
\[
 E_n=\{x:f_n(x)\ge c\varphi(x)\}.
\]
集合 \(E_n\) 单调增加，并且
\(\bigcup_nE_n\supset\{\varphi>0\}\)：若 \(\varphi(x)>0\)，则
\(f_n(x)\uparrow f(x)\ge\varphi(x)>c\varphi(x)\)，故 \(x\) 最终属于
\(E_n\)。于是
\[
 \int f_n\,d\mu\ge c\int_{E_n}\varphi\,d\mu
 \longrightarrow c\int\varphi\,d\mu .
\]
因此 \(L\ge c\int\varphi\,d\mu\)。若某个下方简单函数的积分为无穷，
该式已迫使 \(L=+\infty\)；否则令 \(c\uparrow1\)，得到
\(L\ge\int\varphi\,d\mu\)。再对所有 \(\varphi\le f\) 取上确界，便有
\(L\ge\int f\,d\mu\)。
\end{proof}

若原来处处递增，可以直接定义 \(f=\sup_n f_n\)，其可测性已由
\chapterref{chap:measurable-functions}证明。只知几乎处处收敛时则不能在
任意测度空间上省略极限函数可测的假设：若可测零集有不可测子集 \(Z\)，
取 \(f_n=0\)、\(f=\mathbf1_Z\)，便满足零集外的收敛，但本章并未对这个
不可测的 \(f\) 定义积分。

\begin{corollary}[非负函数的可数可加性]
\label{cor:nonnegative-integral-additivity}
若 \(f,g\ge0\) 可测，则
\[
 \int(f+g)\,d\mu=\int f\,d\mu+\int g\,d\mu.
\]
更一般地，若 \(f_n\ge0\) 可测，则
\[
 \int_X\sum_{n=0}^{\infty}f_n\,d\mu
 =\sum_{n=0}^{\infty}\int_X f_n\,d\mu .
\]
\end{corollary}

\begin{proof}
先用简单函数逼近 \(f,g\)，并对和应用单调收敛定理，得到二项加法。
再令 \(S_N=\sum_{n=0}^{N}f_n\)，则 \(S_N\uparrow\sum_nf_n\)，对
\(S_N\) 应用有限可加性并取极限。
\end{proof}

\section{零积分、水平集估计与密度}

\begin{proposition}[Markov 不等式与零积分判据]
\label{prop:markov-zero-integral}
若 \(f\ge0\) 可测，则对每个 \(t>0\)，
\[
 t\,\mu(\{f\ge t\})\le\int_X f\,d\mu .
\]
并且
\[
 \int_Xf\,d\mu=0\quad\Longleftrightarrow\quad f=0
 \ \text{几乎处处}.
\]
\end{proposition}

\begin{proof}
由 \(t\mathbf1_{\{f\ge t\}}\le f\) 得第一式。若 \(\int f=0\)，则
\(\mu(\{f\ge1/n\})=0\) 对每个 \(n\ge1\) 成立，而
\(\{f>0\}=\bigcup_{n\ge1}\{f\ge1/n\}\)，故 \(f=0\) 几乎处处。反向
若 \(f=0\) a.e.，任意满足 \(0\le\varphi\le f\) 的简单函数也只可能在
零测集上非零，故其积分为零；对这些简单函数取上确界即得反向结论。
\end{proof}

若 \(f\ge0\) 可测，集合函数
\[
 \nu(E)=\int_E f\,d\mu
\]
是 \((X,\mathcal A)\) 上的测度。事实上，对两两不交的 \(E_n\)，有
\(\mathbf1_{\cup E_n}f=\sum_n\mathbf1_{E_n}f\)，再用非负函数的逐项
积分公式。这里 \(f\) 称为 \(\nu\) 关于 \(\mu\) 的密度。反过来哪些测度
能够这样表示，将在后文的 Radon--Nikodym 定理中回答。

\begin{example}[层蛋糕公式的非负版本]
对非负可测函数 \(f\)，有
\[
 f(x)=\int_0^\infty\mathbf1_{\{f(x)>t\}}\,dt.
\]
相应的水平集公式为
\[
 \int_Xf\,d\mu=\int_0^\infty\mu(\{f>t\})\,dt.
\]
其证明安排在习题\ref{ex:simple-integral-level-sets}与
\ref{ex:cavalieri-simple-to-general}中：先处理简单函数，再分别对两个积分
使用单调收敛定理。这条证明适用于任意测度空间，不要求 \(\sigma\)-有限性，
也不需要预先使用乘积测度。
\end{example}

\section*{习题}
\addcontentsline{toc}{section}{习题}

\subsection*{A组　定义与基本运算}
\begin{enumerate}[label=\textbf{\thechapter.\arabic*.}]
\exerciseitem{simple-integral-refinement}
证明命题 \ref{prop:simple-integral-well-defined} 中任意表示公式，并明确处理
某些 \(E_k\) 测度为无穷的情形。

\exerciseitem{simple-integral-level-sets}
设 \(\varphi\ge0\) 为简单函数。证明
\[
 \int\varphi\,d\mu=\int_0^\infty\mu(\{\varphi>t\})\,dt,
\]
右侧按非负函数的扩充值积分理解；当其值有限时，可在包含全部正取值的
有限区间上用通常的一元 Riemann 积分计算。

\exerciseitem{nonnegative-integral-localization}
证明对非负可测函数 \(f\) 和 \(E,F\in\mathcal A\)，若 \(E\cap F=\varnothing\)，
则 \(\int_{E\cup F}f=\int_Ef+\int_Ff\)。

\exerciseitem{nonnegative-integral-ae-invariance}
若 \(f,g\ge0\) 可测且 \(f=g\) 几乎处处，证明
\(\int f=\int g\)。

\end{enumerate}
\subsection*{B组　逼近与收敛}
\begin{enumerate}[label=\textbf{\thechapter.\arabic*.},resume]
\exerciseitem{simple-approximation-alternative}
用水平集 \(\{k2^{-n}\le f<(k+1)2^{-n}\}\) 完整验证定理
\ref{thm:nonnegative-simple-approximation} 的构造。

\exerciseitem{mct-indicator-continuity}
取 \(f_n=\mathbf1_{E_n}\)，从单调收敛定理推出测度关于递增集合列的连续性。
说明这不是循环证明。

\exerciseitem{mct-decreasing-counterexample}
构造 \(f_n\downarrow0\)，但
\(\int f_n\not\to0\) 的非负可测函数列。指出缺少的有限性条件。

\exerciseitem{nonnegative-series-tonelli}
设 \(a_{mn}\ge0\)。用计数测度证明
\[
 \sum_m\sum_na_{mn}=\sum_n\sum_ma_{mn}
\]
允许两边取 \(+\infty\)。

\end{enumerate}
\subsection*{C组　结构与推广}
\begin{enumerate}[label=\textbf{\thechapter.\arabic*.},resume]
\exerciseitem{zero-integral-superlevel}
设 \(f\) 非负可测。证明若存在 \(c>0\) 使
\(\mu(\{f\ge c\})>0\)，则 \(\int f>0\)。给出反例说明
\(f\) 在一个非空稠密集上为正并不推出 \(\int f>0\)。

\exerciseitem{density-measure}
设 \(f\) 非负可测。完整证明 \(\nu(E)=\int_Ef\,d\mu\) 是测度，并证明
\(\nu\ll\mu\)，即 \(\mu(E)=0\Rightarrow\nu(E)=0\)。

\optionalexerciseitem{cavalieri-simple-to-general}
先对简单函数证明层蛋糕公式，再仅用单调收敛思想把它推广到一般非负可测
函数；可以把关于 \(t\) 的积分理解为第 \ref{chap:lebesgue-measure} 章
已构造的 Lebesgue 测度积分。

\projectexerciseitem{integral-from-functional-axioms}
设 \(I\) 是非负可测函数上的泛函，满足对简单函数的积分公式、单调性和
单调收敛性。证明 \(I(f)=\int f\,d\mu\)。这说明本章构造由这些性质唯一
决定。
\end{enumerate}
